\item El 21 de octubre de 2001, Ian Ashpole del Reino Unido logró un récord de altura de $3.35 \text{ km}$ ($11\,000 \text{ pies}$) impulsado por 600 globos llenos con helio. Cada globo lleno tuvo un radio de aproximadamente $0.50 \text{ m}$ y una masa estimada de $0.30 \text{ kg}$. (a) Estime la fuerza de flotación total sobre los 600 globos. (b) Estime la fuerza neta hacia arriba sobre los 600 globos. (c) Ashpole utilizó paracaídas para su retorno a la Tierra después de que los globos empezaron a reventarse a elevadas altitudes y disminuyera la fuerza boyante. ¿Por qué se reventaron los globos?
